\section{The model and the limiting objects}\label{sec:model}

\subsection{Parameters and lattice domains}\label{sec:lattice-domains}
Fix real numbers \(J>0\) and \(U\leq J\) satisfying
\[
 \sinh(2J)=e^{-2U}.
\]
Throughout the proof these parameters remain fixed. Put
\begin{equation}\label{eq:parameters}
 t=(1+e^{4U})^{-1/2},\qquad c=t^{-1},\qquad
 \theta=\arcsin(c/2),\qquad g=\frac{8\theta}{\pi}.
\end{equation}
Then \(1/2\leq t<1\), \(1<c\leq2\), and \(4/3<g\leq4\).
The value \(U=0\) gives \(g=2\); the endpoint
\(J=U=(\log3)/4\) gives \(g=4\). Constants below may depend on the
fixed parameter. We make no claim of uniformity as \(t\uparrow1\).

Let \(D\subset\R^2\) be a bounded Jordan domain. An admissible approximation
is a family of simple nearest-neighbor polygons \(P_\delta\) in
\(\delta\Z^2\) such that their boundaries have parametrizations converging
uniformly to a parametrization of \(\partial D\), and every compact subset
of \(D\) eventually lies inside \(P_\delta\).
Let \(G_\delta=(V_\delta,E_\delta)\) contain all lattice vertices and edges
on or inside the polygon, and let \(B_\delta\) be its boundary-cycle edges.
The weak dual \(G_\delta^\dagger\) has a vertex at each bounded-face center
and an edge across each edge of \(E_\delta\setminus B_\delta\).
We draw all primal and dual edges as straight segments. Boundary
contraction will be used only to describe current weights, not when
drawing clusters.

\subsection{The coupled currents and their height}
We use the spin--edge--height construction of
\cite[Sections~1.1--1.3]{Lis2022Spins}, with its wired/free boundary
specialization made explicit below. Sample spins
\[
 \tau:V_\delta\longrightarrow\{-1,1\},\qquad
 \tau^\dagger:V(G_\delta^\dagger)\longrightarrow\{-1,1\},
\]
with \(\tau=1\) on \(P_\delta\). For each mutually dual pair
\(e=uv\), \(e^\dagger=ff'\), impose
\begin{equation}\label{eq:spin-constraint}
 (\tau(u)-\tau(v))(\tau^\dagger(f)-\tau^\dagger(f'))=0.
\end{equation}
Give an admissible pair weight
\begin{equation}\label{eq:spin-weight}
 t^{N(\tau)+N(\tau^\dagger)},
\end{equation}
where \(N\) counts nearest-neighbor disagreements on the respective
graph. No boundary constraint is imposed on the dual spins.

Conditional on the spins, all edges of \(B_\delta\) are primal open.
For each remaining mutually dual pair, independently of the other pairs,
use the following rule:
\begin{center}
\begin{tabular}{@{}lll@{}}
\toprule
Spin information & Primal/dual open states & Probabilities\\
\midrule
\(\tau^\dagger(f)\ne\tau^\dagger(f')\) & \((1,0)\) & \(1\)\\
\(\tau(u)\ne\tau(v)\) & \((0,1)\) & \(1\)\\
Neither disagreement & \((1,0),(0,1),(0,0)\) & \(1-t,1-t,2t-1\)\\
\bottomrule
\end{tabular}
\end{center}
The constraint makes the first two rows mutually exclusive.
Denote the resulting edge sets by
\(\widehat n_\delta\) and \(\widehat n_\delta^\dagger\).
Their open edges never cross. Both may be closed at a mutually dual pair;
the two currents are not generally complementary.

After contracting the wired boundary and deleting its cycle edges, the
primal odd edges are those crossing dual spin disagreements. The other
primal open edges are even. In the current expansion of
\cite[Section~2, Equations~(2)--(3) and Proposition~1]{Lis2022Boundary},
their law specializes to
\begin{equation}\label{eq:current-weight}
 \P(n_{\rm odd},n_{\rm even})\ \propto\
 2^{k(n_{\rm odd}\cup n_{\rm even})}
 x^{|n_{\rm odd}|}y^{|n_{\rm even}|},
 \qquad
 x=1,\quad y=c-1,
\end{equation}
with \(n_{\rm odd}\) Eulerian and disjoint from \(n_{\rm even}\).
Here \(k\) includes isolated vertices. To check the weight directly,
an odd edge contributes \(t\), an even edge contributes \(1-t\),
and a closed primal edge contributes \(t\), after summing its possible
dual states. The common factor is \(t^{|E_\delta\setminus B_\delta|}\),
leaving \(((1-t)/t)^{|n_{\rm even}|}\).
The primal colors are constant on open components and give
\(2^{k-1}\) choices with the wired color fixed. An admissible Eulerian
odd set determines the dual colors up to a common flip.
These graph-dependent constant factors yield \cref{eq:current-weight}.
The same calculation gives free boundary conditions for the dual current.
In the original parameters,
\[
 x=e^{2U}\sinh(2J)=1,\qquad
 y=e^{2U}\cosh(2J)-1=\sqrt{1+e^{4U}}-1.
\]
It is the spin-and-edge construction above, not just
\cref{eq:current-weight}, that specifies the joint law.

Set \(H_\delta=0\) on \(P_\delta\), and define corner increments by
\begin{equation}\label{eq:height-increments}
 H_\delta(f)-H_\delta(v)=\frac12\tau(v)\tau^\dagger(f)
 \quad\text{when \(v\) is incident to \(f\)}.
\end{equation}
The circulation around an interior primal edge is
\[
 \frac12(\tau(u)-\tau(v))
              (\tau^\dagger(f)-\tau^\dagger(f'))=0.
\]
All boundary primal increments are zero. Since the polygon is simple,
these elementary circuits generate the corner graph's cycle space,
so \cref{eq:height-increments} defines a unique height.
It is integer-valued on primal vertices and half-integer-valued at face
centers. We also use
\[
 K_\delta=2H_\delta.
\]
Thus a nearest corner increment of \(K_\delta\) has magnitude one.
Conversely, the spins are recovered from its residues, up to a common
sign \(\varepsilon\), by
\[
 \tau(v)=\varepsilon(-1)^{K_\delta(v)/2},\qquad
 \tau^\dagger(f)=\varepsilon(-1)^{(K_\delta(f)-1)/2}.
\]
Around a corner-lattice plaquette there are six possible increment
patterns. The two patterns with both opposite diagonals flat have spin
weight \(1\); the other four have weight \(t\). Dividing by the common
factor \(t\) gives the six-vertex weights \((1,1,c)\), with \(c=1/t\);
see also \cite[Section~2.2, Corollary~2.5]{Lis2022Spins}.
The corner lattice is a fixed rotation and dilation of the square
lattice. This local identification, together with the infinite-volume
identification recalled in \cref{sec:spectral}, is what permits the
six-vertex estimates used below.
As a distribution on \(D\), the height is
\begin{equation}\label{eq:height-smearing}
 H_\delta(\varphi)=\delta^2
       \sum_{v\in V_\delta}H_\delta(v)\varphi(v),
 \qquad \varphi\in C_c^\infty(D).
\end{equation}
This is an atomic distribution, which belongs locally to \(H^s\) for
\(s<-1\). Interpolated heights used in intermediate estimates will be
compared with this precise smearing.

Let \(\cC_\delta^{\rm w}\) be the collection of all connected components
of the primal open subgraph, including its boundary component
\(C_{\partial,\delta}\), and let \(\cC_\delta^{\rm f}\) be the collection
of all dual open components. A component includes its vertices and drawn
edges. We retain components at every nesting depth.

\subsection{An auxiliary representation}\label{subsec:auxiliary}
When analyzing one of the two currents, denote it by \(N\), and denote
its exact dual complement by \(W\). This complement is an auxiliary
spin percolation, not the other current of the joint sample.
The spin constant along \(W\) is called \(s\), and the spin constant
along \(N\) is called \(d\). At an ambiguous pair, \(W\) is open with
probability \(t\), whereas the other current is open with probability
\(1-t\).

A \emph{matching wall} means a revealed monochromatic open path or
perimeter in the \(W\)-lattice. Its connections in the revealed region
are retained as boundary connections; separate appearances of the same
connected perimeter are not treated as unrelated wires.
For a free current, the complementary \(W\)-percolation is wired.
The finite cut law in \cref{lem:finite-cuts} explains why the height in
a component behind a matching wall has constant boundary data.
All conditional comparisons below refer to this representation.

The same cut law gives the vocabulary for nesting. A hole of a free
\(N\) component is \emph{odd} when its adjacent \(s\) spin differs from
the adjacent exterior \(s\) spin of that component. Crossing the component
changes \(K_\delta\) by \(2\) times one relative sign in every odd hole,
and by zero in every other hole. Conditional on \((N,s)\), these signs
are independent and fair for distinct free components. After the wired
boundary component is exposed, each residual chamber carries a free
same-parity current; the value of \(K_\delta\) on its component-side
\(W\) perimeter is \(+1\) or \(-1\)
relative to the original boundary height. The edge factor and the jump
identity proving these statements are given in \cref{lem:finite-cuts}.

\subsection{Continuum normalization and two-valued sets}
Let \(h\) be the zero-Dirichlet GFF in \(D\) with covariance
\[
 G_D=2\pi(-\Delta_D)^{-1}.
\]
Thus \(G_D(z,w)=-\log|z-w|+O(1)\) near the diagonal. Set
\begin{equation}\label{eq:normalization}
 \lambda=\frac{\pi}{2},\qquad
 a=\sqrt g\,\lambda,\qquad
 a^2=2\pi\arcsin(c/2).
\end{equation}
In particular \(\pi/\sqrt3<a\leq\pi\), so
\(\lambda<a\leq2\lambda\). A zero-boundary GFF in an open subdomain is
always extended by zero when viewed as a distribution on a larger
domain.

We use the canonical thin two-valued local set
\(\TVS{\alpha}{\beta}(f)\) of a zero-boundary GFF \(f\), with
\(\alpha,\beta>0\) and \(\alpha+\beta\geq2\lambda\). It includes the domain boundary.
The conditional harmonic function on each complementary component is
either \(-\alpha\) or \(\beta\); after subtracting that value, the field in the
component is an independent zero-boundary GFF. These sets are measurable
functions of the field. We use their uniqueness, monotonicity, and
Jordan-hole properties with the hypotheses stated in
\cref{sec:local-sets}; see \cite{ASW2017,AS2018,ALS2017}.
In particular the complementary loops of
\(\TVS{2\lambda}{2\lambda}\) have law \(\CLE_4\).

\begin{definition}[Free cluster procedure]\label{def:free-procedure}
For a zero-boundary GFF \(f\) in a Jordan domain \(O\), define
\(\mathcal F_a(f,O)\) recursively as follows.
\begin{enumerate}
\item Expose \(\TVS{2\lambda}{2\lambda}(f)\). In every complementary
component \(V\), write its label as \(\xi\,2\lambda\),
\(\xi\in\{-1,1\}\), and let \(f_V\) be the zero-boundary remainder.
\item If \(\xi=1\), retain
\(\TVS{2\lambda}{2a-2\lambda}(f_V)\) as one cluster. If \(\xi=-1\),
retain \(\TVS{(2a-2\lambda)}{2\lambda}(f_V)\).
The retained set includes its outer boundary \(\partial V\).
\item In each complementary component of the retained set, subtract its
conditional harmonic value and repeat the procedure with the
zero-boundary remainder.
\end{enumerate}
Keep all retained clusters from all recursive generations. The
\(\CLE_4\) set exposed in the first step is not itself retained.
\end{definition}

If this procedure is invoked above an incoming flat value \(b\) of the
original field, the CLE component has absolute label \(b+\xi2\lambda\)
and the retained split has absolute labels \(b\) and \(b+\xi2a\).
We call the holes with the nonzero jump relative to this incoming value
\(b\) \emph{odd holes}.
The use of this term in lattice explorations is consistent with the
jump decomposition proved below.

Define
\[
 \cC^{\rm f}=\mathcal F_a(h,D).
\]
For the wired collection, first retain
\[
 C_\partial=\TVS{a}{a}(h),
\]
and apply \(\mathcal F_a\) to the zero-boundary remainder in every
complementary component of \(C_\partial\), retaining all resulting
clusters together with \(C_\partial\). Denote that collection by
\(\cC^{\rm w}\). Both constructions use the \emph{same} \(h\), including
the canonical choices of all its local sets. This defines their joint
coupling. The all-generation local finiteness needed below is proved
in \cref{cor:canonical-local-finiteness}, rather than assumed as part of the
definition.

\begin{figure}[!htbp]
\centering
\begin{tikzpicture}[font=\small,>=Latex]
 \node[draw,rounded corners=2pt,align=center,minimum width=30mm,
       minimum height=12mm] (f) at (0,0) {zero-boundary field\\\(f\) in \(O\)};
 \node[draw,rounded corners=2pt,align=center,minimum width=38mm,
       minimum height=12mm] (cle) at (4,0)
       {\(\TVS{2\lambda}{2\lambda}\)\\CLE holes: \(\xi\,2\lambda+f_V\)};
 \node[draw,rounded corners=2pt,align=center,minimum width=38mm,
       minimum height=12mm] (split) at (9,0)
       {retain the split in \(V\)\\relative labels \(0,\xi\,2a\)};
 \draw[->] (f)--(cle);
 \draw[->] (cle)--(split);
 \draw[->] (split.south)--++(0,-.65)-|
       node[pos=.3,below,align=center]
       {subtract the label in each hole; restart independently} (f.south);
\end{tikzpicture}
\caption{One recursive step of the free construction. Only the split
inside a CLE hole is retained as a current cluster. For the wired
construction, the boundary set \(\TVS{a}{a}\) is retained first.}
\label{fig:free-procedure}
\end{figure}

\subsection{Topology and the theorem}
Fix a compact square \(X\) containing \(\overline D\) and all sufficiently
fine polygonal approximations. Such a square exists by admissibility.
Let \(\mathscr H=\mathcal K(X)\) be the space of nonempty compact subsets
of \(X\), with Hausdorff metric \(d_H\), and let
\(Z=\{\{x\}:x\in X\}\subset\mathscr H\) be its singleton subspace.
A collection \(\mathcal K\subset\mathscr H\) is locally finite above
positive diameter if
\[
 \#\{C\in\mathcal K:\diam(C)>\eps\}<\infty
 \qquad\text{for every \(\eps>0\)}.
\]
For such a collection put
\[
 F(\mathcal K)=Z\cup\{C\in\mathcal K:\diam(C)>0\},\qquad
 d_{\rm match}(\mathcal K,\mathcal K')
   =d_H^{\mathscr H}\bigl(F(\mathcal K),F(\mathcal K')\bigr),
\]
where \(d_H^{\mathscr H}\) is Hausdorff distance between compact subsets
of \(\mathscr H\). Local finiteness makes \(F(\mathcal K)\) closed:
any convergent sequence of distinct members has diameters tending to
zero and therefore has a singleton limit. Thus \(d_{\rm match}\) is a
metric after singleton members and multiplicities are ignored.

This metric gives precisely the positive-diameter matching topology.
To make the relation with the usual description explicit, write
\(R_\eps(\mathcal K,\mathcal K')\) for the following two requirements:
\begin{align*}
 &\text{every \(C\in\mathcal K\) with \(\diam(C)>\eps\) has some
 \(C'\in\mathcal K'\) with \(d_H(C,C')<\eps\),}\\
 &\text{and the same statement holds with \(\mathcal K,\mathcal K'\)
 interchanged.}
\end{align*}
The matching need not be one-to-one. We retain the distinguished boundary
component separately, with ordinary Hausdorff topology.

\begin{lemma}[Matching topology and compact containment]
\label{lem:matching-topology}
As relations on locally finite collections,
\[
 \{d_{\rm match}\le\eps/2\}\subset R_\eps
       \subset\{d_{\rm match}\le\eps\}.
\]
Let \(\mathfrak X\) be the compact space of closed subsets
\(F\subset\mathscr H\) containing \(Z\), and write
\(N_t(F)=\#\{C\in F:\diam(C)>t\}\). If \(t_j\downarrow0\) and every
\(M_j\) is finite, then
\[
 \{F\in\mathfrak X:N_{t_j}(F)\le M_j\text{ for every }j\}
\]
is compact and consists of locally finite collections after singleton
members are removed. Consequently, uniform tails for \(N_t\) at each
fixed \(t>0\) imply compact containment in the matching topology.
\end{lemma}

\begin{proof}
For every \(C\in\mathscr H\),
\[
 \frac12\diam(C)\le \dist_{\mathscr H}(C,Z)\le\diam(C).
\]
The lower bound is the triangle inequality, and the upper bound follows
by choosing a singleton in \(C\). If \(d_{\rm match}\le\eps/2\), a member
of diameter greater than \(\eps\) must therefore match a positive-diameter
member of the other family within \(\eps/2\). Conversely, under \(R_\eps\)
each member of diameter at most \(\eps\) is within \(\eps\) of \(Z\), and
each larger member has the required match. This proves the inclusions.

The hyperspace of nonempty compact subsets of a compact metric space is
compact, and the condition \(Z\subset F\) is closed. For each \(t>0\),
\(N_t\) is lower semicontinuous: any finitely many distinct members of
diameter greater than \(t\) persist in disjoint small hyperspace
neighborhoods under Hausdorff convergence. Hence each constraint
\(N_{t_j}\le M_j\) is closed. Their intersection is compact, and choosing
\(j\) with \(t_j<t\) proves finiteness above every \(t>0\). Finally, choose
the \(M_j\) so that the respective tail probabilities are at most
\(\eta2^{-j}\), and use the union bound. This proves compact containment.
\end{proof}

\begin{theorem}[Joint critical-line limit]\label{thm:main}
For every fixed pair $J>0$, $U\le J$ with $\sinh(2J)=e^{-2U}$,
every bounded Jordan domain $D$, and every admissible polygonal
approximation as defined in \cref{sec:lattice-domains},
\[
 \bigl(H_\delta,\cC_\delta^{\rm w},\cC_\delta^{\rm f},
                         C_{\partial,\delta}\bigr)
 \ \xrightarrow[\delta\downarrow0]{\mathrm{law}}\
 \left(\frac{h}{\pi\sqrt g},\cC^{\rm w},\cC^{\rm f},
                         \TVS{a}{a}(h)\right).
\]
The height converges in \(H^s_{\rm loc}(D)\) for every \(s<-1\).
Each cluster collection converges in the matching topology just
defined, and the distinguished component converges in Hausdorff
topology. The continuum collections are locally finite above every
positive diameter. All convergences are joint, and the same GFF
determines both continuum collections.
\end{theorem}

The height topology may equivalently be handled on a countable
exhaustion of \(D\) and a countable sequence of exponents increasing
to \(-1\). Our proof first obtains joint tightness, then identifies
every subsequential limit as the canonical construction of
\cref{thm:main}. This yields the full mesh limit without selecting a
preferred approximation sequence. No claim about cluster area measures,
magnetization fields, all spin-domain-wall curves, or an exact metric
exponent is part of the theorem.

\subsection{Stopping and limiting conventions}\label{subsec:conventions}
All lattice explorations reveal endpoint spins before prescribing an
edge. An accessible chamber keeps the actual boundary connections
created by the explored configuration. At a pinch, its boundary is the
perimeter walk on the side of that chamber; unrelated interiors are
not identified merely because their closures touch.

A \emph{pointed limiting component} is specified by a fixed interior
test point and compact subsets eventually contained in its discrete
component. Several such open components may arise from a single
discrete precursor. Statements about conditional fields in distinct
limiting components therefore include their joint independence after
the appropriate constants are removed. The needed uniform estimates
are established before local-set identification.

Geometric scales, buffers, and diameter thresholds are fixed before
letting \(\delta\downarrow0\). Subsequently a collar width or test scale
may tend to zero, and a finite-count truncation may be removed. Rational
polygons, rational balls, and a fixed countable dense family of interior
points make all uses of these limits simultaneous. Whenever a random
domain is used, its exact cut law specifies the residual boundary class.
An estimate for the residual fields may be invoked conditionally on the
stopping data when it is uniform over that class. Estimates controlling
the production of the cut itself are instead averaged over the exploration,
conditional only on the permitted initial flat data; they are not asserted
after conditioning on a completed stopping history.

\subsection{Proof strategy}\label{subsec:proof-strategy}
The proof first identifies the height field and then recovers the complete
current supports from the geometry coupled to that field.

\Cref{sec:spatial} fixes one current \(N\) and its complementary spin
percolation \(W\). The exact edge factor makes an open monochromatic
\(W\) perimeter a cut with a flat residual law. Ferromagnetic comparison
then supplies circuit estimates in the permitted residual states. The
collar argument controls long passages near a rough explored boundary by
summing bridge failures over the gaps between ordered crossings. It yields
macroscopic count tails and a no-shadowing estimate after a frozen indexed
stop. For a complete prescription of every \(s\) spin and \(W\) edge in a
small region, the pattern comparison compares its probability under allowed
exterior conditions with its plane probability, with a ratio bounded
independently of the number of microscopic prescriptions. A mixed-sign
witness argument then shows that full-spin events with vanishing plane
probability also have vanishing probability uniformly under the allowed flat
laws. This is the one-way contiguity used to transfer screening.

\Cref{sec:spectral} gives Cauchy-mixture representations of the plane
covariance in the four axial and diagonal reflection frames. At nonzero tangential
frequency \(k\), the normal Cauchy width \(s\) has defect
\(D(k,s)=((|k|-s)/(|k|+s))^2\). Fourth angular harmonics in frames differing
by \(45\) degrees cancel pointwise, making this defect integrable against
the mixing measure. For fixed \(0<b<B\), its integral on
\(b\delta\le|k|\le B\delta\) therefore tends to zero, so each rescaled
mixing limit is supported on \(s=|k|\).
Agreement in two independent frames then fixes the remaining spectral weight
in both directions, identifying the Green covariance without a
\(45\)-degree symmetry of the lattice law. In \cref{sec:screening}, the conditional mean of the neutral
Laplacian average \(K_\delta(\Delta\phi)\), given, outside a surrounding square,
both spin fields, the \(W\) states, and computable height increments, tends to zero in \(L^1\). This makes separated
limiting moments harmonic in \cref{sec:gaussian}. Logarithmic bounds then
exclude distributions supported on collision diagonals, and local neutral-pair
mixing identifies the logarithmic collision coefficients, giving Wick's rule.
For plane increment correlations the same argument gives compact-uniform
convergence at every order, which verifies the six-vertex convergence
criterion. Only after it yields the Gaussian plane limit does the free-energy
correspondence fix the coefficient.

\Cref{sec:local-sets} identifies random cuts. The discrete carpet obtained
by stopping each branch at its first odd jump converges to
\(\TVS{2a}{2a}\), which is distinct from a retained cluster of
\cref{def:free-procedure}. Exact finite cuts give independent flat laws in
the residual chambers. A neighborhood argument removes the contributions
of both the revealed mean and the centered remainder near the limiting cut.
Conditional characteristic functions then identify the full zero-extended
GFF remainder, while the piecewise constant labels give the full conditional
mean. This proves the thin local-set law for carpets and the localized
conditional decomposition away from the seed and the initial revealed wall
for cluster stops.
In \cref{sec:local-kernel},
strong convergence of Gaussian projections and cutoff Cameron--Martin
likelihoods give one local cable-test kernel shared by the canonical parent
and a separately stopped component. Testing the canonical forcing event
against the component's zero-boundary law shows that a canonical cluster in
the parent cannot meet both an unsearched stopped component and its boundary.

\Cref{sec:identification} uses this no-crossing result for one inclusion
of the visible partition: odd holes of one canonical cluster must come from
the same discrete precursor. The fair-sign correlation identity excludes a
precursor supplying holes of distinct canonical clusters. Finite witnesses
and the lattice count tails then prove local finiteness across all recursive
generations at once. To recover complete supports, we identify the full
conditional mean of a stopped exploration. An oscillatory GFF test excludes
a connected macroscopic cut where that mean is flat. The intersections of a
connected limiting precursor with the canonical supports and the initial wall
therefore form a countable pairwise disjoint closed cover. Sierpi\'nski's theorem
reduces this cover to one nonempty member. The rough-wall collar estimate
excludes the initial-wall case. The indexed no-shadowing estimate
and the discovery order exclude invisible duplicate precursors. This identifies
every subsequential limit and proves \cref{thm:main}.
